%! Trigonometric Equations and Inequalities — Unit 7, Lesson 7.3 — Exit Ticket (Answer Key)
\documentclass[10pt]{article}
\usepackage{precalculus-article}
\usepackage{precalculus-key}   % pulls in -boxes; do NOT also load -boxes
\newcommand{\TallMath}[1]{$\displaystyle #1\rule[-1.4em]{0pt}{3.2em}$}

\begin{document}

\pageheader{Unit 7, Lesson 7.3}{Exit Ticket --- Answer Key}
\namedateperiod

\vspace{0.10in}

\begin{notesbox}{Exit Ticket --- Answer independently}

\textbf{1.}\quad Solve $\sin\theta = -\dfrac{\sqrt{3}}{2}$ for all $\theta$
in the interval $[0, 2\pi]$.

\vspace{4pt}
Step 1 --- Identify which quadrants have \textbf{negative} sine:
\ans{Quadrants III and IV (sine is negative below the $x$-axis).}

\vspace{4pt}
Step 2 --- Reference angle: $\arcsin\!\left(\dfrac{\sqrt{3}}{2}\right) = $ \ans{$\dfrac{\pi}{3}$}

\vspace{4pt}
Step 3 --- Write both solutions in $[0, 2\pi]$:

\vspace{4pt}
$\theta = $ \ans{$\dfrac{4\pi}{3}$} and $\theta = $ \ans{$\dfrac{5\pi}{3}$}

\begin{teachernote}
Q3 solution: $\pi + \tfrac{\pi}{3} = \tfrac{4\pi}{3}$.
Q4 solution: $2\pi - \tfrac{\pi}{3} = \tfrac{5\pi}{3}$.
Both are in $[0, 2\pi]$.
\end{teachernote}

\vspace{0.16in}
\textbf{2.}\quad Explain in one or two sentences why a trigonometric equation
has \emph{infinitely many} solutions when the domain is all real numbers.

\ans{Trigonometric functions are periodic, so they return to every output value once per period. Adding any integer multiple of the period $2\pi$ to a solution produces another valid solution, giving infinitely many solutions in an unrestricted domain.}

\end{notesbox}

\end{document}
